\chapter{LITERATURE REVIEW}

\section{Wearable Sensor-Based Activity Recognition}
Researchers have increasingly explored wearable sensors as an alternative to
camera-based activity monitoring. Accelerometers and gyroscopes are attractive for
this purpose because they are inexpensive, portable, and able to record movement
without restricting the user to a particular camera view. Surveys by Wang
\textit{et al.} and Lara and Labrador discuss how these sensors are used for activity
recognition, while also pointing out the effects of sensor placement, signal quality,
and differences between users
\cite{wang2019wearable,lara2013survey}.

Kwapisz \textit{et al.} demonstrated that acceleration data can be used to identify
activities using a mobile device, showing the usefulness of time-domain inertial
signals for activity classification \cite{kwapisz2011activity}. Similarly, Bao and
Intille investigated user-annotated acceleration data and showed that sensor
location and the selected activity features influence recognition performance
\cite{bao2004activity}. These findings support the use of two MPU6050 sensors in the
present project, since measurements from two locations can provide more information
about exercise movement and posture than a single sensor.

\section{Inertial Sensors for Exercise Monitoring}
Bulling \textit{et al.} presented a tutorial on human activity recognition using
body-worn inertial sensors and discussed the importance of sensor orientation,
segmentation, and feature representation \cite{bulling2014tutorial}. These issues
are directly relevant to gym-exercise monitoring. A sensor attached at a consistent
body location can record acceleration and angular velocity patterns associated with
the different phases of a squat or bicep curl.

In this project, two MPU6050 devices are used to capture three-axis acceleration and
three-axis gyroscope measurements. Each reading therefore produces a 12-value
feature vector. This arrangement avoids the privacy and viewpoint limitations of
video-based systems while keeping the hardware simple enough for a low-cost gym
assistant prototype.

\section{Evolution of AI in Fitness Monitoring}
Early fitness monitoring systems mainly depended on manual workout records,
pedometers, and simple threshold rules. The development of wearable sensors and
embedded computing enabled continuous collection of physiological and motion data.
Modern AI systems use these measurements to recognize exercises, analyze posture,
count repetitions, and provide feedback during a workout. The AI principles of
perception, learning, and decision making described by Russell and Norvig provide a
general foundation for these intelligent fitness applications
\cite{russell2021aima}.

Recent machine-learning and deep-learning methods can learn increasingly complex
movement patterns from labeled sensor datasets. However, a fitness assistant for a
resource-constrained ESP32 must also consider memory, processing time, power, and
network availability. This project therefore uses a compact KNN model for local
inference while retaining a design that can be extended to more advanced models in
future work.

\section{Machine Learning Methods for Activity Classification}
Activity-recognition systems have used decision trees, support-vector machines,
nearest-neighbor methods, neural networks, and deep-learning models. Geron describes
the importance of preparing features, training a classifier, and evaluating it on
data that was not used during training
\cite{geron2022hands}. McKinney highlights the importance of reliable data loading,
cleaning, analysis, and visualization when building a data-driven system
\cite{mckinney2022python}.

K-nearest neighbors was selected for this project because it is straightforward to
train, works directly with labeled feature vectors, and does not require a large
runtime on the ESP32. The training process standardizes the sensor features and
stores representative examples from each class in program memory. During live use,
the closest examples determine the exercise and posture prediction. This gives the
prototype a model that is easy to inspect and practical to run on the device.

\section{Deep Learning-Based Recognition}
Deep learning has achieved strong results in many perception and activity-recognition
tasks. LeCun, Bengio, and Hinton describe how multilayer neural networks can learn
representations directly from data \cite{lecun2015deep}. Ordonez and Roggen used
convolutional and recurrent neural networks for multimodal wearable activity
recognition, demonstrating the potential of temporal neural models for sensor
sequences \cite{ordonez2016deep}.

Although deep learning is a possible future direction, it is not the primary method
used in the current prototype. Neural-network training generally requires more
data, computation, and memory than the selected compact KNN implementation. KNN
also makes it easier to inspect the relationship between a live sensor vector and
the stored training examples. A future version could compare the embedded KNN model
with a quantized neural network after collecting a larger and more diverse dataset.

\section{Posture Detection and Real-Time Feedback}
Posture correction is an important part of an intelligent fitness assistant because
incorrect form can reduce exercise effectiveness and increase injury risk. A posture
classifier compares live acceleration and angular-velocity patterns with labeled
good- and bad-posture examples. The output can then be converted into immediate
visual or audio feedback.

In the proposed system, the KNN label contains both the exercise and posture state.
The ESP32 displays the predicted exercise, \texttt{good}, \texttt{bad}, or
\texttt{idle} state on an OLED and can activate a buzzer for incorrect posture. The
web dashboard provides the same prediction and confidence for remote monitoring.
This combines recognition and feedback in the same low-cost embedded workflow.

\section{Data Processing and Feature Engineering}
Raw accelerometer and gyroscope readings may contain noise, drift, and fluctuations
caused by sensor placement or user movement. Data preparation is therefore required
before model training. Standardization, consistent sampling, removal of malformed
rows, and validation of class labels improve the reliability of the training data
\cite{mckinney2022python}.

Common motion features include mean acceleration, standard deviation, signal
magnitude, angular velocity, and orientation-related values. The current prototype
uses the twelve direct readings from the two MPU6050 sensors and standardizes them
before calculating KNN distances. This keeps the embedded implementation compact,
while the feature pipeline can later be expanded if additional exercises or users
require greater separation between classes.

\section{Repetition Counting and Movement Segmentation}
Activity recognition identifies what a user is doing, whereas repetition counting
determines how many complete movement cycles have been performed. Ravi \textit{et al.}
and Bao and Intille show that repeated activities can be recognized from recurring
patterns in accelerometer signals \cite{ravi2005activity,bao2004activity}. Signal
processing methods such as smoothing, thresholding, peak detection, and hysteresis
can therefore be used to segment exercise cycles before incrementing a repetition
counter \cite{oppenheim2010discrete}.

In this project, repetition counting is designed as a separate processing stage from
KNN classification. The classifier identifies the exercise and posture, while the
counter uses a filtered motion magnitude and a state transition from an active phase
to a release phase. This separation allows the counting thresholds to be adjusted
for each exercise and prevents a noisy single sample from being interpreted as a
complete repetition.

\section{Embedded and Connected Fitness Systems}
Arduino-based embedded prototyping provides a practical platform for connecting
sensors, displays, actuators, and network services \cite{banzi2022arduino}. The ESP32
extends this approach with wireless communication and real-time task support
\cite{espressif2024esp32}. FastAPI provides a lightweight service layer for receiving
structured sensor batches and serving live status to a browser
\cite{fastapi2024docs}.

The proposed system combines these ideas into one workflow. The ESP32 collects data
from two IMUs, performs local KNN inference in live mode, displays the prediction on
an OLED, and uses a buzzer for posture feedback. In collection mode, labeled samples
are transmitted to the FastAPI service for storage and later model training. The
browser dashboard presents the latest exercise, posture, confidence, and repetition
information to the user.

\section{Research Gap and Proposed Contribution}
Existing activity-recognition research demonstrates the value of wearable inertial
sensors and machine learning, but many solutions focus on general activities,
offline analysis, or computationally larger models. Other systems may identify an
exercise without explicitly providing posture feedback, repetition counting, and a
live user interface in one low-cost prototype.

This project addresses that gap by integrating two MPU6050 sensors, compact KNN
classification, batch voting for stable live predictions, repetition-counting
logic, OLED and buzzer feedback, and a FastAPI web monitor. The current evaluation
focuses on two exercises, bicep curls and squats. The design remains extensible:
additional exercises such as push-ups can be added by collecting labeled samples,
including their classes during training, and regenerating the embedded model.